% ---------------------------------------------------------------------
\addchap{Vorwort}

Fast jede Software, die heute entsteht, wird mit \emph{Git} verwaltet: das Betriebssystem
auf dem Smartphone ebenso wie die App einer Bank oder das Analyseskript einer
Forschungsgruppe. Wer in der Informatik, in Data Science oder im Machine Learning arbeitet,
kommt daran nicht vorbei. In Stelleninseraten wird Git meist gar nicht mehr erwähnt, weil
es schlicht vorausgesetzt wird.

Dieses Handbuch richtet sich an Menschen, die Git \textbf{verstehen} wollen, nicht nur
bedienen. Es setzt keine Vorkenntnisse in Software-Entwicklung voraus. Jeder Fachbegriff
wird beim ersten Auftreten erklärt, und am Ende steht ein Glossar mit den deutschen und
englischen Bezeichnungen.

\section*{Aufbau}

\begin{description}[style=nextline, leftmargin=1.2cm]
  \item[Teil I: Grundlagen] erklärt, welches Problem Versionsverwaltung löst und wie Git
    intern denkt. Dieses \emph{mentale Modell} ist der Schlüssel: Wer es verstanden hat,
    kann sich fast jeden Befehl selbst herleiten.
  \item[Teil II: Arbeiten mit Git] führt durch Einrichtung, täglichen Arbeitsablauf und die
    wichtigsten Befehle, jeweils mit Zweck, Aufbau und Beispiel.
  \item[Teil III: Praxis] zeigt, wie in der Software-Industrie tatsächlich gearbeitet wird:
    gute Commits, Zusammenarbeit, Konflikte, Fehlerbehebung und Sicherheit.
  \item[Anhang] enthält eine Kurzreferenz zum Nachschlagen und das Glossar.
\end{description}

\section*{Konventionen}

\begin{itemize}
  \item Befehle und Dateinamen stehen in \cmd{Schreibmaschinenschrift}.
  \item Befehle, die im Terminal eingegeben werden, stehen in grauen Kästen.
        Alles nach \texttt{\#} ist ein Kommentar und wird nicht mit eingegeben:
\begin{lstlisting}[style=befehle]
git status                  # zeigt den aktuellen Zustand
\end{lstlisting}
  \item Ausgaben des Terminals sind mit einem grauen Balken am Rand markiert:
\begin{lstlisting}[style=ausgabe]
On branch main
nothing to commit, working tree clean
\end{lstlisting}
  \item Platzhalter, die durch eigene Werte ersetzt werden, stehen in spitzen Klammern:
        \cmd{git switch -c <name>} wird zu \cmd{git switch -c ads}.
\end{itemize}

Die farbigen Kästen haben feste Bedeutungen:

\begin{center}
\renewcommand{\arraystretch}{1.3}
\begin{tabularx}{\linewidth}{@{}l>{\raggedright\arraybackslash}X@{}}
  \toprule
  \textbf{Kasten} & \textbf{Inhalt} \\
  \midrule
  {\color{hsluBlau}\textbf{Lernziele / Zusammenfassung}} & was das Kapitel vermittelt \\
  {\color{defFarbe}\textbf{Definition}}                 & präzise Erklärung eines Fachbegriffs \\
  {\color{bspFarbe}\textbf{Aus der Praxis}}             & wie Teams in Unternehmen arbeiten und warum \\
  {\color{gray!65!black}\textbf{Hintergrund}}           & technische Vertiefung, beim ersten Lesen überspringbar \\
  {\color{merkFarbe}\textbf{Erfahrungswert}}            & was man meist erst nach Jahren lernt \\
  {\color{pruefFarbe}\textbf{Typischer Fehler / Achtung}} & häufige Stolpersteine und Gefahren \\
  \bottomrule
\end{tabularx}
\end{center}
